\section{Scope, observations, and host-language assumptions}
\label{sec:scope}
The specification concerns the source snapshot recorded in
\code{sources.sha256}, including the uncommitted rule-structure refactor.
The snapshot is stronger identification than the Git parent alone. The
compiler has five semantic edges:
\[
 \mathsf{CQ}\longrightarrow\mathsf{RA}\longrightarrow\mathsf{IndexedRA}
 \longrightarrow\mathsf{RustFile}[\mathsf{Access}]
 \longrightarrow\mathsf{RustFile}[\mathsf{Pipeline}]
 \longrightarrow\mathsf{RustFile}[\mathsf{Expr}].
\]
The last two edges act on local holes. There is no added nominal IR between
IndexedRA and RustFile. Token emission is not another optimization pass.
Canonical and planned CQ lowering, base and result index selection, and lazy
and materialized storage are separate algorithms at existing edges.

\paragraph{Three distinct claims.}
A local judgment establishes a language's structural well-formedness.
A lowering judgment relates a particular source and target. A preservation
theorem relates their observations. None entails the next without a proof.
An implementation contract can approximate a judgment, and a test can refute
an alleged correspondence, but success does not prove universal preservation.
The assertions called theorems below concern the mathematical rules. Their
application to Rust is conditional on the explicit representation obligations
in Section~\ref{sec:conformance}; we do not claim a verified Rust compiler.

\subsection{Textbook vocabulary: sorts, tuples, and instances}
We use AHV's named/unnamed distinction~\cite[Sections 3.1--3.3,
pp.~30--32]{ahv}. Fix disjoint name sets $\att$ (attributes), $\var$
(variables), and $\mathsf{relname}$ (relation names), and a value domain
$\data$ of constants, disjoint from the name sets. A relation schema $R[U]$ has \emph{sort}
$\sort(R)=U\subseteq\att$ and arity $|U|$. Sorts are finite.
A named tuple over $U$ is a total function $t:U\to\data$;
$t[V]$ denotes its restriction to $V\subseteq U$.
An unnamed tuple of arity $n$ is an element of $\data^n$.
A relation instance $L$ of sort $U$ is a finite set of named tuples over $U$.
A database instance $I$ assigns a relation instance to each relation schema.

The compiler records additional metadata in
$\Sigma(R)=\langle A_0:\tau_0,\ldots,A_{n-1}:\tau_{n-1}\rangle$:
declared order and symbolic Rust types, with
$\sort(R)=\{A_0,\ldots,A_{n-1}\}$. That order relates the named instance
to the ordered Rust input rows. Result names $r$ name further relation
schemas. The implementation's field/function spelling \code{heading}
represents a sort; it introduces no additional semantic object.

\subsection{Valuations and the classical CQ semantics}
\label{sec:valuations}
Following AHV Section~4.2, p.~41, a \emph{valuation} over a finite set
$X\subseteq\var$ is a total map $\nu:X\to\data$. Extend it by
$\nu(a)=a$ for constants and coordinatewise to free tuples, whose entries
are variables or constants. For the range-restricted CQ
$q=\mathsf{ans}(\bar u)\leftarrow R_1(\bar u_1),\ldots,R_n(\bar u_n)$,
\[
q(I)=\{\nu(\bar u)\mid
 \nu:\operatorname{var}(q)\to\data,\quad
 \nu(\bar u_i)\in I(R_i)\text{ for every }i\}.
\]
Here input membership uses ordered tuples in the declared schema order.
This definition gives no variable-assignment or atom-evaluation order.
It is the reference semantics of the classical fragment; Rust computations
and binding clauses require the separate extension below.

\paragraph{Course adaptation: valuation relations.}
Note 1.3~\cite{notes} explicitly calls $\Val_I(C)$ a \emph{valuation relation},
not a named definition from AHV. To make its use of variables as attributes
precise, fix an injective naming map $\iota:\var\to\att$. A valuation
$\nu:X\to\data$ determines the named tuple $\widehat\nu$ by
$\widehat\nu(\iota(x))=\nu(x)$. For a conjunction $C$ of relational atoms,
\[
\Val_I(C)=\{\widehat\nu\mid
 \nu:\vars(C)\to\data,\ I\models C[\nu]\},
\qquad \sort(\Val_I(C))=\iota(\vars(C)).
\]
The course and subsequent rules print $x$ for $\iota(x)$ and suppress the hat
when viewing a valuation as this named tuple. This is an explicit naming
convention; named tuples and valuations have different kinds of domain.
A valuation of $C$ is total on $\vars(C)$ and need not extend to a satisfying
valuation of a larger body. We reserve \emph{partial} for that larger-domain
comparison. The notes' $D$ and our $\data$ denote the same value domain.

\subsection{Named relational operators and the CQ-to-RA connection}
For $\sort(L)=V$ and $\sort(M)=W$, use the named natural join and projection
of AHV Section~4.4, pp.~57--58:
\[
\begin{aligned}
L\bowtie M&=\{t:V\cup W\to\data\mid t[V]\in L,\ t[W]\in M\},\\
\pi_U L&=\{t[U]\mid t\in L\}\quad(U\subseteq V),\\
L\ltimes M&=\pi_V(L\bowtie M),&
\mathbf1&=\{\langle\rangle\}.
\end{aligned}
\]
The semijoin definition is from AHV Section~6.4, p.~128.
For individual tuples $t:V\to\data$ and $s:W\to\data$, write
$t\sim s$ as an abbreviation for $t[V\cap W]=s[V\cap W]$.
Their union is then a function, giving the equivalent expression
$L\bowtie M=\{t\cup s\mid t\in L,s\in M,t\sim s\}$.
Disjoint sorts give Cartesian product. Also
$\pi_\emptyset L=\mathbf1$ iff $L\ne\emptyset$; the empty relation and
the singleton containing the empty tuple remain different.
An injective attribute renaming $m:V\to\att$ acts by
$m(t)(m(A))=t(A)$, and $\delta_m(L)=\{m(t)\mid t\in L\}$,
following AHV's $\delta$ notation (p.~58).

\begin{lemma}[Conjunction as natural join; course derivation]
For relational sub-bodies $C_1,C_2$,
$\Val_I(C_1\land C_2)=\Val_I(C_1)\bowtie\Val_I(C_2)$.
\end{lemma}
\begin{proof}
A satisfying valuation of the conjunction restricts to satisfying valuations
of both sub-bodies. Their named tuples agree on every common attribute,
so their union belongs to the join. Conversely, a compatible pair has a
unique union valuation. Its restrictions satisfy the respective sub-bodies,
so it satisfies the conjunction. These constructions are inverse.
\end{proof}
For distinct-variable atom $A_i=R_i(\bar x_i)$, its positional attribute
renaming $m_i$ therefore satisfies
$\delta_{m_i}(I(R_i))=\Val_I(A_i)$.
Repeated variables and constants require the corresponding position
equalities; a renaming is still injective. Note 1.3 constructs their atom
relations with selection, projection and renaming. The compiler's C/A-rules
implement another construction with generated witness attributes.

For a distinct-variable head $\bar h$ with $H=\set(\bar h)$, define
$\operatorname{out}_{\bar h}(L)=
\{\langle t(h_1),\ldots,t(h_m)\rangle\mid t\in L\}$ for $L$ of sort $H$.
Then, for body $C$,
\[
q(I)=\operatorname{out}_{\bar h}(\pi_H\Val_I(C)).
\]
The projection is a named relation; out renders its tuples in head order.
General heads use the valuation action $\nu(\bar u)$, or the specified Rust
head expression, rather than pretending that projection can repeat columns
or compute values.

\subsection{An explicit interface to Rust}
We do not define all of Rust by a finite DSL expression grammar. Embedded
expressions $e$, patterns $p$, and types $\tau$ are native Syn ASTs. Parsing
an AST does not establish Rust typing, resolution, termination, or purity.
A host evaluation judgment is parameterized by a Rust environment $\rho$,
store $\sigma$, result $v$, and trace $\alpha$:
\[
 \langle e,\rho,\sigma\rangle\Downarrow_H(v,\sigma',\alpha).
\]
There are also panic and divergence observations; absence of a finite
successful judgment does not imply successful termination. The trace may
record I/O, user callbacks, interior mutation, allocation-sensitive events,
and destructor effects. Host iterator creation and polling are separate
operations. A representation relation $v\simeq_\tau a$ relates a concrete
Rust value to an abstract value $a$. Scanned attributes are represented by
borrowed fields; computed attributes are represented by owned values or
lawful clones. This distinction is tracked separately from a sort.

\begin{definition}[Admissible relational host, $\mathcal H_{\rm rel}$]
\label{def:host}
For the reachable expressions and values of an instance:
\begin{enumerate}
\item Every executed expression and pattern match terminates normally, and
patterns used by a Rust \code{let} are irrefutable at their elaborated types.
\item There is a mathematical function $\bar e$ on abstract environments
such that evaluation represents $\bar e(\bar\rho)$, leaves the observable
store unchanged, and emits no observable event. Its result is independent
of representation addresses, aliasing choices, enumeration order, and
irrelevant environment extensions.
\item Generated equality, hashing, cloning, and dropping terminate without
observable effects. Equality is an equivalence relation on abstract values;
equal values hash equally; cloning preserves $\simeq$ and equality. When
sorting is used, comparisons terminate normally without observable effects
and implement a deterministic total order on abstract values consistent
with that equality.
\item Generated Rust is well-typed after name resolution, ownership checking,
and type inference; its access to each logical variable implements the
specified environment interpretation. Allocation failure and resource
exhaustion are outside the mathematical finite-value model.
\end{enumerate}
These are semantic assumptions, not facts established by the current CQ
checker. In particular, \code{Clone + Eq + Hash} trait bounds do not prove
these laws. Syntactically arbitrary Rust remains accepted; this theorem
applies to the stated semantic subdomain of that syntax.
\end{definition}

Under $\mathcal H_{\rm rel}$ define
$\sigma_e L=\{t\in L\mid\bar e(t)=\mathsf{true}\}$ and, for fresh
attributes $N$, $\operatorname{ext}_{N,e}L=\{t\cup\bar e(t)\mid t\in L\}$,
where $\bar e(t):N$. Literal normalization includes the declared Rust column
type: it is not an untyped substitution for every possible host expression.

\subsection{Observational contracts}
For a finite completed cursor $s$, let $\operatorname{Rows}(s)=\set(s)$.
Set-result preservation disregards enumeration order but not values. For a
client request sequence $d$ (build, query, next, drop), define
$\operatorname{Obs}_d(P,I,\sigma)$ to include returned values, panic/divergence,
observable host events, and their stage. Trace preservation requires equality
of these observations for every permitted client. For the abstract cursor
proof below, clients may stop early, but stop polling at their first exhaustion.
This does not assume arbitrary Rust iterators are fused.

The Rust implementations use HashSet iteration, whose order is not prescribed
by the logical language~\cite{hashset}. Under Definition~\ref{def:host}, any
finite enumeration of the represented set is adequate for set results.
Without it, changing enumeration or materialization can change even the final
set. The only inter-route theorem below uses $\operatorname{Rows}$ and
$\mathcal H_{\rm rel}$, not unrestricted trace equivalence.

\begin{counterexample}[Unrestricted route equivalence is false]
\label{cx:effects}
Let $I(R)=\{(10),(20)\}$ and a shared counter initially be zero. Let
\code{tick(x)} increment the counter and return it, and let
\code{latest(n)} test whether the current counter equals $n$:
\begin{lstlisting}
answer(x) :- R(x), let n:i32 = tick(x), if latest(n);
\end{lstlisting}
The lazy route performs tick and its following guard before advancing to the
next input tuple. Both rows survive. The materialized route computes both
extensions before scanning their result for the guard. Exactly the row marked
2 survives, regardless of HashSet order. Moreover, lazy build and cursor
construction call no tick, whereas materialized build calls it twice. A client
that requests one lazy result calls tick once. These are different result and
trace observations. The compiled witnesses in
\code{crates/mini-linq/tests/formal_counterexamples.rs} reproduce both facts.
\end{counterexample}

\paragraph{Consequence.}
Evaluation barriers are necessary dependency constraints; they are not a
proof of cross-route effect preservation. A specification claiming otherwise
would be false. Section~\ref{sec:gaps} records the outstanding language-design
obligation rather than silently weakening an all-Rust correctness claim.
